% ----------------------------------------------------------
\chapter{Introdução}
\label{cap:intr}
% ----------------------------------------------------------

Este documento e seu código-fonte exemplificam a elaboração de trabalhos acadêmicos (trabalhos de conclusão de curso, dissertações e teses) utilizando a classe UnB\TeX, uma extensão da classe abn\TeX2 \cite{abntex2classe} desenvolvida para atender às necessidades da Universidade de Brasília (UnB).

% Definição da nomenclatura que irá para a lista de siglas e abreviações
\nomenclature[A]{UnB}{Universidade de Brasília}

O abn\TeX2, por sua vez, consiste em uma customização da classe \textsf{memoir} destinada à produção de documentos em conformidade com a ABNT NBR 14724, \emph{Informação e documentação -- Trabalhos acadêmicos -- Apresentação}. Informações sobre o projeto estão disponíveis em \url{https://www.abntex.net.br/}.

\nomenclature[A]{ABNT}{Associação Brasileira de Normas Técnicas}

A classe UnB\TeX\ foi desenvolvida com o objetivo de incorporar as atualizações das normas NBR 6023 \cite{NBR6023:2025}, NBR 10520 \cite{NBR10520:2023} e NBR 14724 \cite{NBR14724:2024}, publicadas após a descontinuação das atualizações do abn\TeX2, além de atender aos requisitos específicos dos cursos de graduação e pós-graduação da Universidade de Brasília.

Este documento deve ser utilizado como complemento à documentação do abn\TeX2 \cite{abntex2classe} e da classe \textsf{memoir} \cite{memoir}. Referências adicionais sobre \LaTeX, abn\TeX2 e assuntos correlatos podem ser encontradas em \url{https://github.com/abntex/abntex2/wiki/Referencias}.

%\begin{mdframed}[style=defnSty,innertopmargin=8pt] % azul
\begin{mdframed}[style=plainSty,innertopmargin=8pt] % verde
{\center \textsc{Texto motivador} \par}
\noindent Esperamos que o UnB\TeX\ contribua para a qualidade do seu trabalho, permitindo que sua atenção se concentre principalmente no conteúdo e na contribuição científica.
\end{mdframed}